\documentclass[a4paper,12pt]{article}
\usepackage[T2A]{fontenc}
\usepackage[utf8]{inputenc}
\usepackage[russian]{babel}
\usepackage{amsmath,amssymb,mathtools}
\usepackage{geometry}
\geometry{left=22mm,right=22mm,top=20mm,bottom=21mm}
\usepackage{microtype}
\usepackage{icomma}
\usepackage{tikz}
\usetikzlibrary{arrows.meta,positioning,shapes.geometric,calc}
\usepackage{booktabs,tabularx,array}
\usepackage{enumitem}
\usepackage{fancyhdr}
\usepackage{setspace}
\usepackage{xcolor}
\usepackage[unicode,colorlinks=true,linkcolor=teal!60!black]{hyperref}
\definecolor{accent}{HTML}{087F82}
\definecolor{darkaccent}{HTML}{14565A}
\definecolor{textgray}{HTML}{34393B}
\definecolor{softgray}{HTML}{717A7D}
\definecolor{warmred}{HTML}{A04A48}
\color{textgray}
\onehalfspacing
\setlength{\parindent}{0pt}
\setlength{\parskip}{5pt}
\setlist[itemize]{itemsep=2pt,topsep=3pt,leftmargin=6mm}
\setlist[enumerate]{itemsep=9pt,topsep=5pt,leftmargin=7mm}
\setlength{\headheight}{25pt}
\pagestyle{fancy}
\fancyhf{}
\fancyhead[L]{\small\color{softgray}Проценты: часть, целое и скидки}
\fancyhead[R]{\small\color{softgray}Ксения Васильченко}
\fancyfoot[C]{\small\color{softgray}\thepage}
\renewcommand{\headrulewidth}{0.35pt}
\renewcommand{\headrule}{\hbox to\headwidth{\color{accent}\leaders\hrule height\headrulewidth\hfill}}
\renewcommand{\footrulewidth}{0pt}
\newcommand{\sectionrule}{\par\vspace{-5pt}{\color{accent}\rule{\linewidth}{0.5pt}}\vspace{5pt}}
\newcommand{\ctitle}[1]{\section*{\color{darkaccent}#1}\sectionrule}
\newcommand{\subt}[1]{\vspace{5pt}{\color{darkaccent}\bfseries #1}\par}
\newcommand{\keyidea}[1]{\textcolor{darkaccent}{\textbf{Главная мысль.}} #1}
\begin{document}
\begin{titlepage}
\thispagestyle{empty}
\begin{center}
\vspace*{21mm}
{\color{darkaccent}\Large\textbf{МАТЕМАТИКА}}\par
\vspace{22mm}
{\fontsize{26}{31}\selectfont Чек-лист}\par
\vspace{10mm}
{\color{accent}\fontsize{21}{26}\selectfont\bfseries Проценты: часть, целое\par и задачи на скидки}\par
\vspace{17mm}
{\large Задачи на проценты: от схемы к решению}\par
\vfill
{\large Ксения Васильченко}\par
\vspace{10mm}
{\normalsize 12 ноября 2026 г.}\par
\vfill
{\small Лёвин Артём Александрович}\par
{\small эксклюзивно для Ксении Васильченко}\par
\vspace{19mm}
\begin{tikzpicture}[x=1cm,y=1cm]
\draw[accent,line width=.8pt] (-5.1,0) .. controls (-3.7,0.75) and (-2.0,-.45) .. (-.55,0)
  .. controls (1.0,.6) and (3.25,-.65) .. (5.1,0);
\fill[accent] (-.55,0) circle (1.25pt);
\end{tikzpicture}
\end{center}
\end{titlepage}

\ctitle{1. Три величины в задаче на проценты}
В каждой задаче сначала выясняйте, что принято за \textbf{100\%}. Эта величина называется
\textbf{целым}. Затем определяйте, какая \textbf{часть} рассматривается и \textbf{какой процент}
целого она составляет.

\begin{center}
\renewcommand{\arraystretch}{1.25}
\begin{tabularx}{\linewidth}{@{}p{30mm}Xl@{}}
\toprule
Величина & Что означает & Обозначение\\
\midrule
Целое & Всё количество; ему соответствует $100\%$ & $C$\\
Часть & Количество внутри целого & $A$\\
Процент & Сколько сотых долей целого составляет часть & $p\%$\\
\bottomrule
\end{tabularx}
\end{center}
Число $p$ в записи $p\%$ обязательно переводим в \textbf{долю}:
\[
\boxed{d=\frac{p}{100}},\qquad
6\%=\frac{6}{100}=0{,}06,\quad
30\%=\frac{30}{100}=0{,}3,\quad
85\%=\frac{85}{100}=0{,}85.
\]

\subt{Три формулы, которые нужно знать}
\[
\begin{array}{ll}
\textbf{Найти часть:}   & \displaystyle A=C\cdot\frac{p}{100},\\[10pt]
\textbf{Найти целое:}  & \displaystyle C=A:\frac{p}{100}\quad (p\ne0),\\[10pt]
\textbf{Найти процент:}& \displaystyle p=\frac{A}{C}\cdot100\quad (C\ne0).
\end{array}
\]

\subt{«Треугольник процентов»}
\begin{center}
\begin{tikzpicture}[scale=1, every node/.style={align=center}]
\draw[darkaccent,line width=1.25pt] (-2.6,-.9)--(0,2.2)--(2.6,-.9)--cycle;
\draw[darkaccent,line width=1pt] (-1.45,.62)--(1.45,.62);
\draw[darkaccent,line width=1pt] (0,.62)--(0,-.9);
\node[font=\small\bfseries] at (0,1.25) {ЧАСТЬ $A$};
\node[font=\small\bfseries] at (-1.25,-.17) {ДОЛЯ $\dfrac{p}{100}$};
\node[font=\small\bfseries] at (1.25,-.17) {ЦЕЛОЕ $C$};
\end{tikzpicture}
\end{center}
\keyidea{Сверху стоит часть. Снизу --- доля и целое. Чтобы получить верхнюю величину, перемножьте нижние; чтобы найти нижнюю, разделите верхнюю на известную нижнюю.}

\clearpage
\ctitle{2. Как правильно читать условие и строить схему}
Для задач с несколькими процентами используйте \textbf{схему-отрезок}. Над отрезком записывайте
количества в рублях, килограммах или штуках; под ним --- проценты. Весь отрезок равен $100\%$.

\subt{Пример 1. Расходы на ремонт}
На ремонт потратили 275~тыс.~руб. Из них $30\%$ --- оплата рабочим,
$50\%$ --- стройматериалы, остальное --- сантехника. Сколько денег потратили на сантехнику?
\begin{center}
\begin{tikzpicture}[x=1cm,y=1cm]
\draw[darkaccent,line width=1pt] (0,0) rectangle (13.2,1.0);
\draw[darkaccent,line width=1pt] (3.96,0)--(3.96,1);
\draw[darkaccent,line width=1pt] (10.56,0)--(10.56,1);
\node[font=\small,align=center] at (1.98,.5) {рабочие};
\node[font=\small,align=center] at (7.26,.5) {материалы};
\node[font=\small,align=center] at (11.88,.5) {сантехника};
\node[font=\small] at (1.98,-.34) {$30\%$};
\node[font=\small] at (7.26,-.34) {$50\%$};
\node[font=\small] at (11.88,-.34) {$20\%$};
\draw[{Stealth[length=2mm]}-{Stealth[length=2mm]},darkaccent] (0,1.3)--(13.2,1.3);
\node[fill=white,inner sep=3pt,font=\small] at (6.6,1.3) {275 тыс. руб. = $100\%$};
\end{tikzpicture}
\end{center}
\textbf{Шаг 1.} Найдём неизвестный процент: $100-30-50=20\%$.

\textbf{Шаг 2.} Найдём $20\%$ от всей суммы:
\[
275\cdot\frac{20}{100}=55\ \text{(тыс. руб.)}.
\]
\textcolor{darkaccent}{\textbf{Ответ:}} 55 тыс. руб.

\subt{Пример 2. Сахар в клубнике}
В 12{,}5~кг клубники содержится $6\%$ сахара. Сколько килограммов сахара?
\[
A=12{,}5\cdot\frac{6}{100}=0{,}75\ \text{кг}.
\]
Здесь $12{,}5$~кг --- целое, $6\%$ --- известный процент, а масса сахара --- искомая часть.

\subt{Пример 3. Распределение бюджета}
На семейный отдых выделили 320~тыс.~руб. На билеты потратили $25\%$ всей суммы,
на проживание --- $45\%$ всей суммы. Остальное --- на экскурсии и питание.
\[
100\%-25\%-45\%=30\%,\qquad
320\cdot\frac{30}{100}=96\ \text{(тыс. руб.)}.
\]
\textcolor{warmred}{\textbf{Важно:}} оба процента относятся ко всей первоначальной сумме.
Поэтому вычитаем $25\%$ и $45\%$ из $100\%$, сохраняя одно и то же целое.

\clearpage
\ctitle{3. Как найти целое и процент}
\subt{Пример 4. Найти целое по известной части}
В спортивной секции 24 человека --- это $30\%$ всех участников. Сколько всего человек?
\[
C=24:\frac{30}{100}=24:0{,}3=80\ \text{(человек)}.
\]
\textbf{Проверка:} $80\cdot0{,}3=24$.\par
\textcolor{darkaccent}{\textbf{Ответ:}} 80 человек.

\textbf{Приём для деления на десятичную дробь:}
\[
24:0{,}3=240:3=80.
\]
Чтобы делитель стал целым числом, умножьте и делимое, и делитель на 10.

\subt{Пример 5. Найти процент по части и целому}
В питомнике 120 саженцев, из них 42 --- берёзы. Какой процент саженцев составляют берёзы?
\[
\frac{42}{120}=0{,}35,\qquad
p=0{,}35\cdot100=35.
\]
\textcolor{darkaccent}{\textbf{Ответ:}} $35\%$.

\subt{Пример 6. Проверка на новых числах}
На выставке 150 книг, из них 54 --- научно-популярные. Их доля равна
\[
\frac{54}{150}=0{,}36,\qquad p=36.
\]
\textcolor{darkaccent}{\textbf{Ответ:}} $36\%$ книг.

\subt{Как быстро выбрать действие}
\begin{itemize}
\item \textbf{Дано целое и процент} $\longrightarrow$ найдите часть: умножьте.
\item \textbf{Дано часть и процент} $\longrightarrow$ найдите целое: разделите на долю.
\item \textbf{Дано часть и целое} $\longrightarrow$ найдите процент: разделите часть на целое и умножьте на 100.
\end{itemize}
\keyidea{Числовая проверка помогает обнаружить ошибку: если нашли $30\%$ от количества, результат должен составлять примерно треть этого количества.}

\clearpage
\ctitle{4. Скидки: от первоначальной цены к новой и обратно}
Скидка уменьшает первоначальную цену. Если скидка составляет $p\%$,
то после неё остаётся $(100-p)\%$ первоначальной стоимости.

\subt{Пример 7. Известна цена до скидки}
Куртка стоила 1200~руб. Магазин снизил цену на $15\%$. Сколько стоит куртка теперь?
\[
100\%-15\%=85\%,\qquad
1200\cdot\frac{85}{100}=1020\ \text{руб.}
\]
\begin{center}
\begin{tikzpicture}[x=1cm,y=1cm]
\draw[darkaccent,line width=1pt] (0,0) rectangle (13.2,.85);
\draw[darkaccent,line width=1pt] (11.22,0)--(11.22,.85);
\node[font=\small] at (5.61,.43) {новая цена $=85\%$};
\node[font=\small] at (12.21,.43) {скидка};
\node[font=\small] at (5.61,-.32) {1020 руб.};
\node[font=\small] at (12.21,-.32) {$15\%$};
\draw[{Stealth[length=2mm]}-{Stealth[length=2mm]},darkaccent] (0,1.18)--(13.2,1.18);
\node[fill=white,inner sep=3pt,font=\small] at (6.6,1.18) {1200 руб. = $100\%$};
\end{tikzpicture}
\end{center}
\textbf{Проверка:} $15\%$ от 1200 руб. --- 180 руб.; $1200-180=1020$ руб.

\subt{Пример 8. Известна цена после скидки}
Спортивная сумка после скидки $20\%$ стоит 840~руб. Сколько она стоила сначала?
\[
100\%-20\%=80\%,\qquad
C=840:\frac{80}{100}=840:0{,}8=1050\ \text{руб.}
\]
\textbf{Проверка:} $1050\cdot0{,}8=840$ руб.

\subt{Пример 9. Другой пример обратной скидки}
После скидки $20\%$ конструктор стоит 1040~руб. Тогда первоначальная цена:
\[
C=1040:0{,}8=1300\ \text{руб.}
\]
\textcolor{warmred}{\textbf{Частая ошибка:}} прибавлять к новой цене $20\%$ от неё.
Скидка считалась от \emph{первоначальной} цены. Поэтому сначала найдите,
какой процент первоначальной цены составляет стоимость после скидки.

\clearpage
\thispagestyle{fancy}
\begin{center}
{\Large\color{darkaccent}\bfseries Алгоритм решения задачи на проценты}\par
\vspace{7mm}
\end{center}
\begin{center}
\begin{tikzpicture}[
  >=Stealth,
  start/.style={ellipse,draw=darkaccent,line width=.9pt,minimum width=7cm,minimum height=1.1cm,align=center,font=\small},
  action/.style={rounded corners=2pt,rectangle,draw=darkaccent,line width=.9pt,align=center,font=\small,minimum height=1.3cm,text width=8.4cm,inner sep=6pt},
  choice/.style={diamond,draw=darkaccent,line width=.9pt,aspect=3.0,align=center,inner sep=2pt,font=\small,text width=4cm},
  result/.style={rounded corners=2pt,rectangle,draw=accent,line width=1pt,align=center,font=\small,minimum height=1.7cm,text width=4.25cm,inner sep=4pt},
  arr/.style={->,darkaccent,line width=.9pt}
]
\node[start] (start) at (0,0) {Прочитайте условие};
\node[action] (all) at (0,-2.0) {Определите \textbf{целое}: что равно $100\%$?\\Запишите данные и единицы.};
\node[choice] (which) at (0,-4.65) {Что нужно найти?};
\node[result] (part) at (-5.3,-8.1) {\textbf{ЧАСТЬ}\\[4pt] $A=C\cdot\dfrac{p}{100}$\\[3pt] Умножьте на долю.};
\node[result] (whole) at (0,-8.1) {\textbf{ЦЕЛОЕ}\\[4pt] $C=A:\dfrac{p}{100}$\\[3pt] Разделите на долю.};
\node[result] (pct) at (5.3,-8.1) {\textbf{ПРОЦЕНТ}\\[4pt] $p=\dfrac{A}{C}\cdot100$\\[3pt] Найдите отношение.};
\node[action] (check) at (0,-11.45) {Проверьте единицы и смысл ответа.\\Подставьте результат в условие.};
\node[start] (end) at (0,-13.45) {Запишите ответ};
\draw[arr] (start)--(all);
\draw[arr] (all)--(which);
\draw[arr] (which.south west) -| node[pos=.30,above,sloped,font=\small] {часть} (part.north);
\draw[arr] (which.south)--node[right,font=\small,near start] {целое} (whole.north);
\draw[arr] (which.south east) -| node[pos=.30,above,sloped,font=\small] {процент} (pct.north);
\draw[darkaccent,line width=.9pt] (part.south)--(-5.3,-10.1)--(5.3,-10.1)--(pct.south);
\draw[darkaccent,line width=.9pt] (whole.south)--(0,-10.1);
\draw[arr] (0,-10.1)--(check.north);
\draw[arr] (check)--(end);
\end{tikzpicture}
\end{center}
\clearpage

\ctitle{5. Типичные ошибки и итоговая памятка}
\begingroup\setstretch{1.08}\setlength{\parskip}{2pt}\small
\subt{Ошибка 1. Не выделено целое}
В задаче о клубнике $100\%$ --- это все 12{,}5~кг ягод, а $6\%$ --- доля сахара.
Прежде чем считать, назовите величину, соответствующую $100\%$.

\subt{Ошибка 2. Забыли разделить процент на 100}
\[
12{,}5\cdot6=75\quad\text{--- неверно},\qquad
12{,}5\cdot\frac{6}{100}=0{,}75\quad\text{--- верно}.
\]

\subt{Ошибка 3. Перепутали часть и целое}
Если 24 человека --- это $30\%$ секции, то $24:0{,}3=80$ человек.
Результат больше 24: ведь мы нашли всю секцию, а исходные 24 человека были её частью.

\subt{Ошибка 4. Неверно определили процент остатка}
В задаче о бюджете $25\%$ и $45\%$ всей суммы потрачены на две разные цели.
Остаток составляет $100-25-45=30\%$.

\subt{Ошибка 5. Неверно определили базу скидки}
Если вещь после скидки $20\%$ стоит 840~руб., то 840~руб. соответствуют $80\%$
\emph{начальной} цены. Делим $840$ на $0{,}8$.

\subt{Ошибка 6. Потеряли скобки в вычислении}
\[
15-3-2=10,\qquad 15-(3-2)=14.
\]
Скобки показывают, какие действия выполняются первыми.

\vspace{5pt}
\subt{Короткая памятка}
\begin{center}
\renewcommand{\arraystretch}{1.35}
\begin{tabularx}{\linewidth}{@{}p{56mm}X@{}}
\toprule
Сигнал в условии & Что делать\\
\midrule
«Найдите $p\%$ от числа» & Целое умножить на $p/100$.\\
«$A$ --- это $p\%$ целого» & Часть разделить на $p/100$.\\
«Какой процент составляет $A$ от $C$?» & $A$ разделить на $C$, умножить на 100.\\
«Осталось после расходов / скидки» & Сначала найти оставшийся процент.\\
«Цена после скидки известна» & Принять её за $100\%-p\%$ исходной цены.\\
\bottomrule
\end{tabularx}
\end{center}
\endgroup

\clearpage
\ctitle{Тренировочный блок}
\textcolor{softgray}{\small Путь А --- 10 задач. Решайте по порядку. Для каждой задачи сначала определите, что принято за $100\%$, и выберите схему или треугольник.}\par
\vspace{4pt}
\textbf{Базовый уровень}
\begin{enumerate}[label=\textbf{\arabic*.}]
\item Найдите $14\%$ от 250.
\item В 12{,}5~кг фруктов содержится $8\%$ сахара. Сколько килограммов сахара в этих фруктах?
\item В кружке 27 ребят --- это $30\%$ всех участников. Сколько всего ребят в кружке?
\item В питомнике 180 саженцев, из них 45 --- яблони. Какой процент саженцев составляют яблони?
\end{enumerate}
\vspace{3pt}
\textbf{Средний уровень}
\begin{enumerate}[label=\textbf{\arabic*.},start=5]
\item На ремонт потратили 36\,000~руб. Рабочим заплатили $25\%$ суммы, на материалы ушло $40\%$, остальное --- на сантехнику. Сколько рублей потратили на сантехнику?
\item Куртка стоила 2400~руб. После скидки $15\%$ сколько рублей нужно заплатить за куртку?
\item Сумка после скидки $20\%$ стала стоить 720~руб. Сколько она стоила до скидки?
\item На выставке 160 книг, из них 64 --- научно-популярные. Какой процент книг составляют научно-популярные?
\end{enumerate}
\vspace{3pt}
\textbf{Уровень выше среднего}
\begin{enumerate}[label=\textbf{\arabic*.},start=9]
\item На поездку выделили 520~тыс.~руб. На билеты потратили $25\%$ всей суммы, на гостиницу --- $45\%$ всей суммы. Сколько тысяч рублей осталось на остальные расходы?
\item После скидки $15\%$ конструктор стал стоить 1700~руб. Найдите его первоначальную цену и величину скидки в рублях.
\end{enumerate}

\clearpage
\ctitle{Ответы к тренировочному блоку}
\textcolor{softgray}{\small Сверяйте ответы после самостоятельного решения. В ответах сохранены единицы измерения.}
\vspace{6pt}
\begin{center}
\renewcommand{\arraystretch}{1.65}
\begin{tabularx}{.88\linewidth}{@{}>{\centering\arraybackslash}p{16mm}X@{}}
\toprule
№ & Ответ\\
\midrule
1 & $35$\\
2 & $1$ кг\\
3 & $90$ человек\\
4 & $25\%$\\
5 & $12\,600$ руб.\\
6 & $2040$ руб.\\
7 & $900$ руб.\\
8 & $40\%$\\
9 & $156$ тыс. руб.\\
10 & Первоначальная цена --- $2000$ руб.; скидка --- $300$ руб.\\
\bottomrule
\end{tabularx}
\end{center}
\vfill
\begin{center}
\color{softgray}\small Лёвин Артём Александрович эксклюзивно для Ксении Васильченко
\end{center}
\end{document}
